\documentclass[11pt]{article}
\usepackage{fontspec}
\usepackage{booktabs}
\usepackage{iftex}
\setmainfont{FreeSerif}[RawFeature={-ccmp}]
\newcommand{\EvidenceMode}{synthetic}
\newcommand{\RealCount}{zero}
\newcommand{\FrozenCovariate}{kappa E}
\newcommand{\FrozenFamily}{logistic}
\begin{document}
\section*{Regime-adaptive rank fusion of W1 and message passing}
\textbf{Evidence mode.} The available continuation checkout contained no
recoverable real-network acquisition manifest or pilot branch. The analysis is
therefore reported as evidence from synthetic graphs only; \RealCount{} real
networks were available in this checkout.

\textbf{Method.} W1 and two-iteration Karrer--Newman--Zdeborová message passing
are reused as members. Their ascending vertex ranks are combined as
$F_c=(1-c)r(W1)+cr(MP2)$. The selection objective uses the independent
development sample and never uses evaluation labels.

\textbf{Covariates and selection.} The primary covariate is the mean edge
probability over existing edges. Challengers include expected degree,
escape probability, matrix and non-backtracking spectral radii, message
passing giant-component mass, and the configuration-model branching factor.
The recorded development candidate uses \FrozenCovariate{} with a
\FrozenFamily{} response.

\textbf{Limitations.} The missing pilot artifacts prevent a faithful
reconstruction of the requested v1 comparison, and the available acquisition
environment did not yield a real network. Results must not be interpreted as
real-network evidence.
\end{document}
